% Last updated: 2025/04/22

\subsection{Descriptive comparison}
Figure \ref{fig:scatter} shows a simple scatterplot comparison between a surgeon's current treatment intensity and their medical school HRR intensity. As evidenced by the graph, there seems to be little difference between the mover and stayer groups. This descriptive result will be quantified with later estimates. Interestingly, about 67\% of surgeons in my sample practice less intensely than the average intensity they were exposed to in medical school. This proportion is similar among movers and stayers, though slightly higher in the stayer group (70\%). This below-average intensity is likely due to the fact that newer surgeons with less experience may be more selective about when to operate due to lack of experience. 

\begin{figure}
     \centering
     \begin{subfigure}[b]{0.45\textwidth}
         \centering
         \includegraphics[width=\textwidth]{intensity_corr_movers.png}
         \label{fig:scatter-movers}
     \end{subfigure}
     \hfill
     \begin{subfigure}[b]{0.45\textwidth}
         \centering
         \includegraphics[width=\textwidth]{intensity_corr_stayers.png}
         \label{fig:scatter-stayers}
     \end{subfigure}
    \caption{Correlation between med school HRR intensity and current intensity}
    \label{fig:scatter}
\end{figure}

\subsection{Effect of level of medical school intensity}
In Table \ref{tab:res-main-short}, I present the main estimates of equation \ref{eq:main}. I estimate the equation in the full sample of all surgeons, as well as two subgroups consisting of only general surgeons and only specialist surgeons. I consider a specialist surgeon to be any surgeon whose specialty is not general surgery. General surgeons undergo residency training that is often the basis for other surgical subspecialties. I expect to see a larger impact of medical school within generalists, since they tend to undergo less GME compared to surgical specialists. 

\setlength{\LTpost}{0mm}
\LTcapwidth=\textwidth
\begin{ThreePartTable}
\begin{TableNotes}[para, flushleft]
    \textit{Notes.} Clustered standard errors are based on the physician's current HRR Model 1 includes indicators for all surgical subspecialties, and model 3 includes indicators for surgical subspecialties other than vascular surgery. All models include physician gender, patient characteristics, and indicators for missing patient characteristics, as shown in Table \ref{tab:summ}. All models also include current HRR fixed effects and year fixed effects. $* p < 0.10, ** p < 0.05, *** p < 0.01$.
\end{TableNotes}
\begin{longtable}{lccc}
    %\centering
    \label{tab:res-main-short} \\
    \caption{Effect of medical school intensity on physician own intensity} \\
    %\begin{tabular}{lccc}
        \toprule
         & (1) & (2) & (3) \\
         & All surgeons & Generalists & Specialists \\
        \midrule
        Mover $\times$ med school HRR intensity & -0.017 & 0.134 & -0.018 \\ 
         & (0.106) & (0.093) & (0.117) \\ 
        Med school HRR intensity & -0.223 & -0.203 & -0.237 \\ 
         & (0.278) & (0.222) & (0.313) \\ 
        \midrule
        Observations & 6760 & 723 & 6037 \\ 
        $R^2$ & 0.233 & 0.725 & 0.231 \\ 
        Mean physician intensity & 0.210 & 0.139 & 0.219 \\ 
        \bottomrule
        \insertTableNotes
    %\end{tabular}
\end{longtable}
\end{ThreePartTable}


The effect of medical school HRR intensity is statistically insignificant across the board. Examining the magnitude of effect for the full sample reveals that medical school seems to have no meaningful effect on how a physician practices now. A 1 percentage point increase in medical school HRR intensity decreases a surgeon's current treatment intensity by about 0.017 percentage points compared to their peers in the same destination HRR. This corresponds to about a 0.08\% decrease in treatment intensity. Although this is not a tight null, the 95\% confidence interval suggests that the effect of one percentage point increase in medical school HRR intensity ranges between a 1\% decrease and a 1\% increase in current physician intensity. The entire range is not economically meaningful. The effect estimated within the specialist subsample is similarly both statistically and economically insignificant. 

The magnitude of the effect within the generalist subsample is larger, as expected, but still statistically insignificant. A positive effect here indicates that general surgeons who attended medical school in more intense areas practice more intensely compared to their current peers. The effect size is still small, however. A 1 percentage point increase in medical school HRR intensity increases a general surgeon's current practice intensity by 0.134 percentage points, which corresponds to about a 1\% increase in intensity. The 95\% confidence interval suggests that the effect is ranges between a -0.35\% decrease and a 2.3\% increase.

Estimates of the semi-parametric specification in equation \ref{eq:spline} reveal that the effect of medical school intensity is not much different at the tails, although estimates are fairly noisy. In Figures \ref{fig:spline}, I present predicted physician intensity using only the cubic spline term. The magnitudes of predicted physician intensity when predicting using the mean of all other covariates is similar, but much noisier. The unadjusted effect of medical school intensity is small, and medical school intensity does not have a differential effect between movers and stayers. 
\begin{figure}
     \centering
     \begin{subfigure}[b]{0.3\textwidth}
         \centering
         \includegraphics[width=\textwidth]{spline-full.png}
         \label{fig:spline-full}
     \end{subfigure}
     \hfill
     \begin{subfigure}[b]{0.3\textwidth}
         \centering
         \includegraphics[width=\textwidth]{spline-general.png}
         \label{fig:spline-general}
     \end{subfigure}
     \hfill
     \begin{subfigure}[b]{0.3\textwidth}
         \centering
         \includegraphics[width=\textwidth]{spline-specialist.png}
         \label{fig:spline-specialist}
     \end{subfigure}
    \caption{Predicted physician intensity without controls}
    \label{fig:spline}
\end{figure}

\subsection{Effect of difference in intensity}
Results from specifications utilizing the change in intensity from medical school HRR to destination HRR tell a slightly different story compared to examining the level of medical school intensity alone. The raw level of medical school intensity does not seem to have an effect, which may be confounded by some surgeons moving to more intense areas, while other surgeons move to less intense areas. Estimates of equation \ref{eq:delta} allows for greater comparison using the differences between the intensity in the training market compared to the practice market. 

Estimated effects are presented in Table \ref{tab:res-delta-short}. Results in the full sample and specialist sample are statistically significant at the 0.10 level, but not at the 0.05 level. The negative sign here indicates that an increase in the difference between a surgeon's current practice HRR and their medical school decreases their current treatment intensity relative to their peers in the same current HRR. That is, the more intense a surgeon's current practice HRR is compared to their medical school HRR, the more they tend to practice less intensely compared to their peers. Specifically, the increasing the difference by 1 percentage point decreases a surgeon's current practice intensity by 2.5 percentage points. This corresponds to a 12\% decrease. The 95\% confidence interval suggests the effect is between a 26\% decrease and a 2\% increase. It's possible that surgeons who train in less intense areas adjust their treatment intensity to accommodate norms in their current practice market, but only to the extent of the support of their medical school HRR intensity distribution. Further work, perhaps with more granular data, is needed to determine how the distribution of the medical school HRR impacts current practice intensity in a way beyond just the mean. 

\setlength{\LTpost}{0mm}
\LTcapwidth=\textwidth
\begin{ThreePartTable}
\begin{TableNotes}[para, flushleft]
    \textit{Notes.} Clustered standard errors are based on the physician's current HRR Model 1 includes indicators for all surgical subspecialties, and model 3 includes indicators for surgical subspecialties other than vascular surgery. All models include physician gender, patient characteristics, and indicators for missing patient characteristics, as shown in Table \ref{tab:summ}. All models also include current HRR fixed effects and year fixed effects. $* p < 0.10, ** p < 0.05, *** p < 0.01$.
\end{TableNotes}
\begin{longtable}{lccc}
    %\centering
    \label{tab:res-delta-short} \\
    \caption{Effect of change in intensity on physician own intensity} \\
    %\begin{tabular}{lccc}
        \toprule
         & (1) & (2) & (3) \\
         & All surgeons & Generalists & Specialists \\
        \midrule
        Mover $\times$ change in intensity & -2.544* & -0.412 & -2.919* \\ 
         & (1.503) & (0.323) & (1.701) \\ 
        Med school HRR intensity & -2.783* & -0.505 & -3.171* \\ 
         & (1.664) & (0.389) & (1.881) \\ 
        \midrule
        Observations & 6760 & 723 & 6037 \\ 
        $R^2$ & 0.241 & 0.723 & 0.241 \\ 
        Mean physician intensity & 0.210 & 0.139 & 0.219 \\
        \bottomrule
        \insertTableNotes
    %\end{tabular}
\end{longtable}
\end{ThreePartTable}


These findings suggest that medical school intensity does not have a statistically or economically significant effect on a physician's current practice. Overall, this provides support for the story where medical school provides a foundational knowledge base, but practical knowledge about patient types and when to operate is developed and refined in residency and fellowship, contrary to popular belief. It's possible that medical knowledge is taught differently in medical schools, but the information that a surgeon retains from medical school about treatment intensity seems to be somewhat standardized. 